\documentclass[10pt,a4paper]{amsart}
\usepackage[margin=19mm]{geometry}
\usepackage{amsmath,amssymb,mathtools}
\usepackage[scr=rsfs,cal=euler]{mathalfa}
\usepackage{xfrac}
\usepackage[unicode,hidelinks,bookmarks=false]{hyperref}
\newcommand{\ie}{i.e.\ }
\newcommand{\cf}{cf.\ }
\newcommand{\resp}{respectively\ }
\newcommand{\Subsection}{\subsection}
\providecommand{\nobreakdash}{\nobreak\hbox{-}}
\setlength{\emergencystretch}{2em}
\multlinegap0pt
\usepackage{xcolor}
\usepackage{amssymb}
\usepackage{tikz}
\usepackage{xfrac}
\usepackage{float}
\usepackage{stackengine}
\usepackage[shortlabels]{enumitem}
\usepackage{mathtools}
\newtheorem{prop}{Proposition}
\newtheorem{lemma}{Lemma}
\title{Some cusp-transitive hyperbolic 4-manifolds}
\author{Edoardo Rizzi}
\date{Source: arXiv:2402.13209v1 (CC BY 4.0)}
\begin{document}
\maketitle

\noindent\emph{Source and licence.} Some cusp-transitive hyperbolic 4-manifolds. Version arXiv:2402.13209v1 (CC BY 4.0), distributed by arXiv under the Creative Commons Attribution 4.0 International licence, http://creativecommons.org/licenses/by/4.0/, as recorded in the arXiv licence metadata for this exact version and shown on the abstract page https://arxiv.org/abs/2402.13209v1. Full licence notice: \texttt{licenses/arxiv-2402.13209-CC-BY-4.0.txt}.\par\smallskip

\noindent\textit{Editorial scope.} Counted unit: the article's lemma corner (source0.tex lines 1622--1635). The article's complete original proof of it is delivered with it (source0.tex lines 1636--1646, one \texttt{proof} environment of 11 lines). No mathematics is written, repaired or re-derived: the statement and the proof are the article's own lines, in the article's own notation, and no retained line is substituted or re-flowed.  Retained uncounted as the source-local context the proof needs: lines 756--759; lines 1584--1591.  Nothing was omitted from any retained span, so the package carries no omission record.  No retained line contains a citation, so the wrapper carries no bibliography and no bibliography entry.  No retained line contains a figure environment, an image-inclusion command or drawing code, so no image is delivered and the package declares no asset dependency.  Editorial formatting only, in addition: an AMS \texttt{amsart} preamble that restates the article's own declarations and macros used by the retained lines, namely the environments \texttt{prop}, \texttt{lemma} on separate counters, together with the standard packages those lines need.  The retained environments print in this document as Proposition 1, Lemma 1 in order of appearance, whereas the article prints the same items with the numbering of its own full document.  The complete native source of the article as distributed in the arXiv e-print for this exact version is bundled unchanged as \texttt{sources/154873/source0.tex}.
\begin{prop}
\label{angolo}
    If two facets of type $i$ and $j$ of $P_n$ meet, with $i\ne j$, then the dihedral angle between them is the same of the one between the facets $\textbf{i}$ and $\textbf{j}$ of $P_0$. In particular the facets $\textbf{i}$ and $\textbf{j}$ of $P_0$ meet.
\end{prop}

\begin{lemma}
\label{tipi}
    A facet of $R$ is:
    \begin{itemize}
        \item either the image through $p$ of a facet of type $7$,
        \item or the image through $p$ of a union of facets of type $3$.
    \end{itemize}
\end{lemma}

\begin{lemma}
\label{corner}
    A corner of $R$ is:
    \begin{itemize}
        \item either of type $(3,3)$, and in this case it is the image through $p$ of a union of some type-$(3,3)$ ridges;
        \item either of type $(7,7)$, and in this case it is:
        \begin{itemize}
            \item either the image through $p$ of a type-$(7,7)$ ridge of $P$;
            \item or the image through $p$ of a type-$(7,i)$ ridge of $P$, with $i=1,6$;
        \end{itemize}
        
        \item or of type $(3,7)$, and in this case it is the image through $p$ of a type-$(3,7)$ ridge of $P$.
    \end{itemize}
\end{lemma}

\begin{proof}
    Since the facets of $R$ are of type $3$ or $7$, there are three kinds of corners in $R$: type $(3,3)$, $(7,7)$, and $(3,7)$.

    By Proposition \ref{angolo}, if a facet of type $3$ and a facet of type $1$ or $6$ of $P$ meet, the dihedral angle between them is $\frac{\pi}{2}$. Hence the image through $p$ of a type-$(3,i)$ ridge of $P$, with $i=1,6$, is contained in the relative interior of a type-$3$ facet. Hence the union of the type-$(3,3)$ corners of $R$ is the image through $p$ of the union of the type-$(3,3)$ ridges of $P$.

    The image through $p$ of a type-$(3,3)$ ridge is contained in a corner, hence every type-$(3,3)$ corner is the image through $p$ of the union of some type-$(3,3)$ ridges.
    
    %Given $A$ and $B$ two facets of $P$ that we identify with the gluing. Let $S_A$ and $S_B$ be the set of facets of type $3$ of $P$ that they respectively meet. Let $M,N\in S_A$

    Since by Lemma \ref{tipi} the image through $p$ of a facet of type $7$ of $P$ is a facet of type $7$ of $R$, the image of a type-$(7,7)$, or type-$(3,7)$, ridge is a corner of $R$. Moreover,  the image of a type-$(7,i)$ ridge, with $i=1,6$, is a type-$(7,7)$ corner.
\end{proof}

\end{document}
